\documentclass[12pt,a4paper]{article}

% XeLaTeX: шрифти й мови
\usepackage{fontspec}
\usepackage{polyglossia}
\setmainlanguage{ukrainian}
\setotherlanguage{english}
\setmainfont{Times New Roman}
\newfontfamily\cyrillicfonttt{Courier New}

% Макет сторінки
\usepackage{geometry}
\geometry{top=2.5cm,bottom=2.5cm,left=2.5cm,right=2.5cm}

% Графіка й таблиці
\usepackage{graphicx}
\usepackage{booktabs}
\usepackage{longtable}

% Математика
\usepackage{amsmath, amsfonts, amssymb}

% Код
\usepackage{listings}
\usepackage{xcolor}
\lstdefinestyle{pythonstyle}{
    language=Python,
    basicstyle=\ttfamily\small,
    keywordstyle=\color{blue},
    commentstyle=\color{gray},
    stringstyle=\color{orange!60!black},
    numberstyle=\tiny\color{gray},
    numbers=left,
    numbersep=8pt,
    frame=single,
    breaklines=true,
    tabsize=4,
    showstringspaces=false
}

% Гіперпосилання
\usepackage{hyperref}

% Підписи підграфіків
\usepackage{caption}
\usepackage{subcaption}
\begin{document}

%-------------------- Title page --------------------
\begin{titlepage}
    \centering
    {\large\bfseries Національний технічний університет України\\
    ``Київський політехнічний інститут імені Ігоря Сікорського'' \par}
    \vspace{2cm}
    {\Large Розрахунково-графічна робота\par}
    \vspace{1cm}
    {\LARGE\bfseries Порівняння алгоритмів ройового інтелекту (ACO, GA, Tabu Search)\\
    для задачі маршрутизації транспортних засобів з часовими вікнами\par}
    \vfill
    \begin{flushright}
    Виконав: студент гр.Фі-21\\
    Грунда Я.І.\\[0.5cm]
    Перевірив: ас. Железняков Д.В.
    \end{flushright}
    \vfill
    {\large Київ --- \the\year\par}
\end{titlepage}

\tableofcontents
\newpage

%-------------------- Актуальність -------------------
\section{Актуальність теми}
Транспортні витрати становлять до 60\,\% операційного бюджету більшості логістичних компаній. 
Раціональне планування маршрутів із урахуванням часових вікон доставки 
(\emph{Vehicle Routing Problem with Time Windows}, VRPTW) дає змогу 
скоротити пробіг, кількість залучених автомобілів і, відповідно, 
зменшити витрати пального та викиди \(\mathrm{CO_2}\). 
Серед метаевристик, що практично застосовуються, часто розглядають 
алгоритми ройового інтелекту (Ant Colony Optimization, Genetic Algorithm) 
та траєкторні методи (Tabu Search). 
Тому доцільно оцінити продуктивність цих алгоритмів на вибірці тестових інстансів VRPTW.

\section{Мета та завдання роботи}
\textbf{Мета}: порівняти якість ACO, GA та Tabu~Search на
еталонних інстансах Solomon VRPTW.

\textbf{Завдання}:
\begin{enumerate}
  \item Вивчити математичну постановку VRPTW.
  \item Реалізувати версії ACO, GA, Tabu Search з однаковою
        цільовою функцією.
  \item Розробити інструмент автоматичного запуску та логування історії cost.
  \item Провести серію експериментів з однаковим тайм‑лімітом.
  \item Порівняти мінімальні, середні та максимальні показники cost і distance,
        зробити висновки.
\end{enumerate}


\section{Постановка задачі VRPTW}

Нехай задано орієнтований граф \(G=(V,E)\), де  
\(V=\{0,1,\dots,n\}\): вершина \(0\)~— депо, а \(1..n\) — клієнти.  
Дуга \((i,j)\in E\) характеризується відстанню \(d_{ij}\)  
(дорівнює часу проїзду).  
Для кожного клієнта \(i\ (i\ge 1)\) задано

\begin{itemize}
  \item попит \(q_i\);
  \item часовий інтервал~\([e_i,l_i]\), у межах якого має
        розпочатися обслуговування;
  \item тривалість сервісу \(s_i\).
\end{itemize}

У депо знаходиться флот із \(K_{\max}\) однакових вантажівок
ємністю \(Q\) та власним часовим вікном \([e_0,l_0]\).
Необхідно побудувати набір маршрутів  
\(\mathcal R=\{R_1,\dots,R_K\}\) (\(K\le K_{\max}\)),  
кожний з яких починається й закінчується в депо, так щоб

\begin{enumerate}
  \item кожен клієнт відвіданий рівно один раз;
  \item \emph{вантажні} та \emph{часові} обмеження виконуються
        (або штрафуються) для кожного маршруту;
  \item мінімізується наступна цільова функція:
\end{enumerate}

\[
\min\;
D
+ \lambda_{\text{late}}\sum_{i=1}^{n} L_i
+ \lambda_{\text{over}}\sum_{k=1}^{K} O_k
+ \lambda_{\text{depot}}\sum_{k=1}^{K} D_k
+ \lambda_{\text{veh}}\,
      \max\bigl(0,\,K-K_{\max}\bigr),
\]

де
\begin{description}
  \item[\(D\)] — сума всіх відстаней по маршрутах;
  \item[\(L_i\)] — \textbf{запізнення} на клієнті \(i\):
        \(L_i=\max\bigl(0,\;T_i - l_i\bigr)\),  
        \(T_i\) — фактичний час початку сервісу;
  \item[\(O_k\)] — \textbf{перевантаження} маршруту \(k\):
        \(O_k=\max\bigl(0,\;\sum_{i\in R_k} q_i - Q\bigr)\);
  \item[\(D_k\)] — \textbf{порушення вікна депо} для маршруту \(k\):  
        \(D_k=\max\bigl(0,\;t^{\text{back}}_k - l_0\bigr)\),  
        де \(t^{\text{back}}_k\) — час повернення у депо;
  \item[\(\lambda_{\ast}\)] — великі штрафні коефіцієнти,
        що збігаються з \texttt{lambda\_late}, \texttt{lambda\_over},
        \texttt{lambda\_depot}, \texttt{lambda\_veh}
        у функції \texttt{evaluate\_solution}.
\end{description}

Таким чином, якщо всі обмеження виконуються,
штрафові суми \(L_i,\;O_k,\;D_k\) дорівнюють нулю,
і функція зводиться до мінімізації маршрутової відстані \(D\)
та кількості використаних машин.

\subsection{Параметри штрафів}
\begin{table}[h]
\centering\begin{tabular}{lcc}
\toprule
Коефіцієнт & Значення & Коментар\\
\midrule
$\lambda_\text{over}$ & 1\,000 & заборона перевантажень авто\\
$\lambda_\text{late}$ & 1\,000 & заборона запізнень до клієнтів\\
$\lambda_\text{depot}$ & 1\,000 & заборона пізнього повернення до депо\\
$\lambda_\text{veh}$  & 10\,000 & жорстке обмеження кількості авто\\
\bottomrule
\end{tabular}
\caption{Налаштування цільової функції}
\end{table}

\section{Опис алгоритму Ant Colony Optimization}
\subsection{Загальна ідея ACO}
Ant Colony Optimization (ACO) належить до метаевристик ройового
інтелекту і моделює колективну поведінку колонії мурах, що
шляхом відкладання \emph{феромону} знаходять найкоротші стежки до
джерела їжі.  На рівні алгоритму це відбувається так:

\begin{enumerate}
  \item Кожен агент (``мураха'') послідовно будує розв’язок,
        переходячи зі вершини \(i\) у вершину \(j\) із вірогідністю,
        пропорційною \(\tau_{ij}^{\alpha}\,\eta_{ij}^{\beta}\),  
        де \(\tau_{ij}\)~— інтенсивність феромону,  
        \(\eta_{ij}\)~— статична евристика  
        (для VRPTW~— \(1/d_{ij}\)).
  \item Після побудови всіх розв’язків:
        \begin{itemize}
          \item феромон \emph{випаровується}  
                \(\tau_{ij}\leftarrow(1-\rho)\tau_{ij}\);
          \item кращі розв’язки \emph{накопичують} феромон  
                \(\tau_{ij}\leftarrow\tau_{ij}+ \Delta\tau_{ij}\),
                підсилюючи ребра, якими варто ходити.
        \end{itemize}
  \item Процедура повторюється, доки не буде вичерпано ліміт
        ітерацій або часу.  У процесі мурахи дедалі частіше
        використовують ребра, що ведуть до глобально кращих
        траєкторій.
\end{enumerate}

Таким чином ACO комбінує \textbf{різноманітність} (рандомізовані
ходи) та \textbf{інтенсифікацію} (підсилення кращих ребер) і
поступово концентрує пошук поблизу найякісніших рішень.

%-------------------------------------------------
\subsection{Деталі реалізованої модифікації}
\subsubsection{Структури даних}
\begin{itemize}
  \item \(\tau\in\mathbb R^{(n+1)\times(n+1)}\) — матриця феромону,
        ініціалізована константою~\(\tau_0\).
  \item \(\eta_{ij}=1/(d_{ij}+10^{-6})\) — евристична привабливість.
  \item \texttt{history}: список кортежів
        \((\text{iter},\,t,\,\text{best\_cost},\,\text{best\_dist})\)
        для графіків збіжності.
\end{itemize}

\subsubsection{Параметри алгоритму}
\begin{table}[h]
\centering
\begin{tabular}{lcl}
\toprule
Параметр & Значення & Значення за замовчуванням \\
\midrule
Кількість мурах $n_{\text{ants}}$ & $10\dots 50$ & 50 \\
$\alpha$ (вага феромону)          & 1–2 & 1.0 \\
$\beta$ (вага евристики)          & 2–4 & 3.0 \\
$\rho$ (випаровування)            & 0.1–0.3 & 0.1 \\
$q_0$ (ймовірність жадібного кроку) & 0.3–0.7 & 0.6 \\
$\tau_0$ (старт. феромон)          & $10^{-4}\dots 10^{-2}$ & $10^{-3}$ \\
\texttt{elitist}                  & True/False & True \\
\bottomrule
\end{tabular}
\caption{Ключові параметри реалізації ACO}
\end{table}

\subsubsection{Крок побудови туру}
\begin{enumerate}
  \item Старт у депо \(0\); множина невідвіданих клієнтів~\(U\).
  \item Для поточного вузла \(i\) обчислити
        \(\text{desir}_j=\tau_{ij}^{\alpha}\eta_{ij}^{\beta}\) для
        всіх \(j\in U\).
  \item З імовірністю \(q_0\) вибрати
        \(j=\arg\max_{j\in U}\text{desir}_j\);  
        інакше провести рулетку по \(\text{desir}_j\).
  \item Додати \(j\) до туру, виконати \emph{локальне} зменшення
        феромону:  
        \(\tau_{ij}\leftarrow(1-\rho)\tau_{ij}+\rho\tau_0\).
  \item Повторювати, поки \(U=\varnothing\);  
        на завершення виконати local-update для ребра
        останній $\to$ депо.
\end{enumerate}

\subsubsection{Глобальне оновлення феромону}
Після кожної ітерації:
\[
\tau_{ij}\leftarrow(1-\rho)\tau_{ij} 
      + \sum_{(i,j)\in R_\star}\frac{1}{\text{cost}(R_\star)},
\]
де $R_\star$ — глобально найкращий тур (режим \texttt{elitist}).
У не-elitist-режимі сума береться по всіх мурах поточної ітерації.

\subsubsection{Критерій зупинки}
Алгоритм працює, доки одночасно виконується
\[
\text{elapsed}<\text{time\_limit}\quad\land\quad
\text{iter}<\text{max\_iter}.
\]
Обидва параметри задаються користувачем; кожен може бути~`None`.

\subsubsection{Запис історії}
Після кожної ітерації у \texttt{history} додається запис
\((k,\;t,\;\text{best\_cost},\;\text{best\_dist})\),  
що дозволяє будувати графіки «best-cost / distance vs time»
і порівнювати збіжність з іншими алгоритмами.

\subsubsection{Повернене значення}
Метод \texttt{run()} віддає пару
\[
(\text{best\_cost},\;\text{best\_tour}),
\]
де \texttt{best\_tour} — giant-tour (перестановка клієнтів),
а вартість \texttt{best\_cost} враховує відстань та всі штрафи,
визначені у функції \texttt{evaluate\_solution}.

\section{Опис алгоритму Genetic Algorithm (GA)}
\subsection{Загальна ідея GA}
Генетичний алгоритм (GA) — це популяційна метаевристика,
що імітує еволюційні процеси:  
\emph{відбір найкращих}, \emph{кросовер} (рекомбінація ДНК)  
та \emph{мутацію}.  
У контексті VRPTW:

\begin{enumerate}
  \item Кожен індивід кодує \textbf{giant-tour} —
        перестановку клієнтів \(1..n\); депо $0$ додається
        під час декодування.
  \item Популяція розміром \(P\) оцінюється цільовою функцією
        (відстань + штрафи за запізнення, перевантаження тощо).
  \item На кожному поколінні
        \begin{enumerate}
          \item обираємо батьків оператором турнірного відбору;
          \item створюємо потомка кросовером OX (Order Crossover);
          \item випадково мутуємо потімка, міняючи місцями 2–3 гени;
          \item еліта (найкращий індивід) копіюється без змін.
        \end{enumerate}
  \item Після  завершення ліміту часу або 
        заданої кількості поколінь повертається глобальний
        найкращий розв’язок.
\end{enumerate}

Таким чином GA забезпечує \textbf{різноманітність}
за рахунок мутацій і рекомбінацій,  
а \textbf{інтенсифікацію} — через відбір кращих
та елітизм.

%-------------------------------------------------
\subsection{Деталі реалізованої модифікації}
\subsubsection{Структури даних}
\begin{itemize}
  \item \texttt{Individual} = (tour, cost),  
        де \texttt{tour} — список клієнтів,  
        \texttt{cost} — значення цільової функції.
  \item \texttt{pop} — список із \(P\) особин.
  \item \texttt{best} — глобально найкращий індивід.
  \item \texttt{history} --- кортежі
        \((\text{покоління},\,t,\,
        \text{best\_cost},\,\text{best\_dist})\).
\end{itemize}

\subsubsection{Параметри алгоритму}
\begin{table}[h]
\centering
\begin{tabular}{lcl}
\toprule
Параметр & Діапазон / сенс & Значення за замовчуванням \\
\midrule
$P$ (\texttt{pop\_size})  & 100–500  & 300 \\
Ставка кросоверу $p_\text{cx}$ & 0.7–0.95 & 0.9 \\
Ставка мутації  $p_\text{mut}$ & 0.05–0.3 & 0.2 \\
Розмір турніру $k$            & 2–5 & 3 \\
Еліта                         & 1 індивід & +1 \\
\bottomrule
\end{tabular}
\caption{Ключові параметри GA-реалізації}
\end{table}

\subsubsection{Оператори}
\paragraph{Відбір}  
$k$-арний турнір без повернення:  
обираємо $k$ випадкових індивідів,  
беремо того, чия \texttt{cost} мінімальна.

\paragraph{Кросовер} — \emph{Order Crossover (OX).}  
Зберігає відносний порядок клієнтів:
\begin{enumerate}
  \item випадковий сегмент [\(a,b\)) копіюється з~батька~1;
  \item решту позицій послідовно заповнюють гени батька~2,
        зберігаючи їхню відносну чергу.
\end{enumerate}

\paragraph{Мутація} — \emph{swap.}  
З імовірністю $p_\text{mut}$ міняємо місцями 2–3 довільні позиції.

\paragraph{Елітизм}  
Найкращий індивід популяції копіюється у нове покоління без змін.

\subsubsection{Критерій зупинки}
Цикл триває, доки виконується
\[
\text{elapsed} < \text{time\_limit}
\quad\land\quad
\text{generation} < \text{max\_gen}.
\]
Будь-який з параметрів може бути \texttt{None}.

\subsubsection{Запис історії}
Наприкінці кожного покоління:
\[
\texttt{history}\; \mathrel{+\!=}\;
(g,\;t,\;\text{best\_cost},\;\text{best\_dist}),
\]
що дозволяє порівнювати криві збіжності з ACO та Tabu Search.

\subsubsection{Повернене значення}
\texttt{run()} повертає
\[
(\text{best\_cost},\;\text{best\_tour}),
\]
де \texttt{best\_tour} — giant-tour із мінімальним \texttt{cost}
за всю еволюцію.

\section{Опис алгоритму Tabu Search}
\subsection{Загальна ідея Tabu Search}
Tabu Search (TS) — \emph{траєкторна} метаевристика: вона оперує
одним робочим розв’язком і крок-за-кроком рухається його околом.
Основні принципи:

\begin{itemize}
  \item \textbf{Найкращий доступний сусід.}  
        На ітерації формується підмножина можливих рухів
        (swap/relocate). Вибирають той, чия вартість мінімальна,
        навіть якщо вона \emph{гірша} за поточну.  
        Так пошук може «перестрибувати» локальні мінімуми.
  \item \textbf{Короткострокова памʼять (tabu-list).}  
        Застосований рух заносять у список табу довжини
        $\textit{tenure}$ і тимчасово забороняють —  
        це запобігає негайному поверненню та циклам.
  \item \textbf{Аспірація.}  
        Якщо табу-рух генерує розв’язок, який
        \emph{покращує глобальний best}, заборона ігнорується.
        Отже, «хороші винятки» не блокуються.
  \item \textbf{Довгострокова еволюція.}  
        Чергуючи керовані погіршення і заборони,
        алгоритм поступово досліджує нові області простору
        й концентрується навколо перспективних зон.
\end{itemize}

Такий баланс \emph{інтенсифікації} (пошук найкращого сусіда)
і \emph{диверсифікації} (табуювання та аспірація)
дозволяє TS ефективно обходити локальні пастки
та знаходити якісні рішення VRPTW.

\begin{enumerate}
  \item Початкова точка — випадковий giant-tour.
  \item На кожній ітерації генерується підмножина
        сусідів (swap/relocate), оцінюється їхня вартість,
        і вибирається ліпший \emph{допустимий} кандидат.
  \item Виконаний рух заноситься в табу-список довжини
        \texttt{tenure}; прострочені записи автоматично
        видаляються.
  \item Алгоритм завершується після тайм-ліміту
        або коли досягнуто \texttt{max\_iter}.
\end{enumerate}

\subsection{Деталі реалізованої модифікації}
\subsubsection{Структури даних}
\begin{itemize}
  \item \textbf{Move} = \((i,j,\text{type})\),  
        де \(\text{type}\in\{\text{``swap''},\,\text{``reloc''}\}\).
  \item \texttt{tabu} — \textbf{черга} записів \((\text{move},\,\text{expiry})\).
  \item \texttt{history} — \((\text{iter},\,t,\text{best\_cost},\text{best\_dist})\).
\end{itemize}

\subsubsection{Параметри алгоритму}
\begin{table}[h]
\centering
\begin{tabular}{lcl}
\toprule
Параметр & Діапазон / сенс & Значення за замовчуванням \\
\midrule
Довжина tabu~\textbf{tenure}            & 10–40 & 20 \\
Розмір підмножини о́круги~\textbf{neigh\_size} & 100–500 & 200 \\
Частка \texttt{swap} у сусідах~\textbf{swap\_rate} & 0.3–0.7 & 0.5 \\
\bottomrule
\end{tabular}
\caption{Ключові параметри Tabu Search}
\end{table}

\subsubsection{Сусідство}
\begin{itemize}
  \item \textbf{Swap}\((i,j)\): міняє місцями позиції \(i\) і \(j\) у giant-tour.
  \item \textbf{Relocate}\((i\!\to\!j)\): виймає клієнта з позиції \(i\)
        і вставляє перед \(j\) (або після, щоб зберегти індекси).
\end{itemize}
На кожній ітерації випадково відбирається
\texttt{neigh\_size} рухів згідно зі \texttt{swap\_rate}.

\subsubsection{Алгоритм однієї ітерації}
\begin{enumerate}
  \item Для кожного кандидата \(c\) обчислюємо  
        \(\text{cost}(c)=\)~\texttt{evaluate\_solution}.
  \item Відкидаємо рух \emph{tabu}, якщо він не дає кращий
        global~best (аспірація).
  \item Обираємо кандидата з мінімальним \(\text{cost}\);
        робимо його новим ``current''.
  \item Якщо \(\text{cost}<\text{best\_cost}\) —
        оновлюємо глобальний best.
  \item Додаємо рух у \texttt{tabu} зі строком
        дії \(\text{iter}+\texttt{tenure}\).
\end{enumerate}

\subsubsection{Запис історії}
Після кожної ітерації:
\[
\texttt{history}\mathrel{+\!=}
(it,\;t,\;\text{best\_cost},\;\text{best\_dist}).
\]

\subsubsection{Критерій зупинки}
Пошук триває, доки
\[
\text{elapsed}<\text{time\_limit}
\quad\land\quad
\text{iter}<\text{max\_iter}.
\]

\subsubsection{Повернене значення}
Метод \texttt{run()} повертає
\[
(\text{best\_cost},\;\text{best\_tour}),
\]
де \texttt{best\_tour} — giant-tour з мінімальним штрафним cost,
знайдений за весь час роботи.

\section{Експериментальна частина}
\subsection{Налаштування тестів}
Тайм-ліміт кожного запуску~— 60 с.  
Для інстансу \textbf{C101} виконано 10 запусків,  
для решти~— по 5 запусків з різними початковими seed.

\vspace{0.5em}
\begin{table}[ht]
\centering
\caption{Параметри запуску кожного алгоритму}
\label{tab:setup}
\begin{tabular}{@{}lccc@{}}
\toprule
\textbf{Параметр} & \textbf{ACO} & \textbf{GA} & \textbf{Tabu Search}\\
\midrule
Популяція \,/\, мурахи            & $n_{\text{ants}} = 25$ & $P = 25$ & --- \\
Вага феромону $\alpha$            & 1.0 & --- & --- \\
Вага евристики $\beta$            & 1.5 & --- & --- \\
Коеф. випаровування $\rho$        & 0.1 & --- & --- \\
$ q_0 $ (жадібність)              & 0.3 & --- & --- \\
Початковий феромон $\tau_0$       & $1\times10^{-3}$ & --- & --- \\
Elitist-режим                     & True & --- & --- \\
\addlinespace
Турнір $k$                        & --- & 5  & --- \\
Ставка кросоверу $p_\text{cx}$    & --- & 0.9 & --- \\
Ставка мутації  $p_\text{mut}$    & --- & 0.8 & --- \\
\addlinespace
Tenure табу-списку                & --- & --- & 15 \\
Розмір околу \textit{neigh\_size} & --- & --- & 250 \\
\bottomrule
\end{tabular}
\end{table}

\subsection{Опис інстансів Solomon}

Інстанси Solomon VRPTW поділяються на три класи:

\begin{itemize}
  \item \textbf{C-серія (Clustered)} — клієнти згруповані у декілька щільних
        кластерів; корисні для перевірки алгоритму на
        «багато коротких рейсів усередині кластерів».
  \item \textbf{R-серія (Random)} — координати клієнтів рівномірно
        розкидані по всій області; дають довгі радіальні маршрути.
  \item \textbf{RC-серія (Random‐Clustered)} — змішаний тип,
        деякі точки утворюють кластери, решта — розкидані випадково.
\end{itemize}

У кожній серії є підсерії \textbf{1} та \textbf{2}:
\(\text{C}1, \text{R}1, \text{RC}1\) мають \emph{вузькі} часові вікна,  
через що потрібна більша кількість машин;  
серії~\(\text{C}2, \text{R}2, \text{RC}2\) мають
\emph{ширші} вікна й зазвичай дозволяють менше транспортних засобів.

\subsection{Результати}
\begin{longtable}{lcccccc}
\toprule
Алгоритм & min cost & avg cost & max cost & min dist & avg dist & max dist\\ 
\midrule
ACO  & 2491404.42  & 2534961.06  & 2584992.51 & 1889.96 & 1949.43 & 2048.73\\
GA   & 3171845.28  & 4155713.68 & 5140755.60 & 3480.26 & 3672.15 & 3955.61\\
Tabu & 586228.99  & 1208464.63 & 1435677.21 & 2655.69  & 2986.31 & 3410.03\\
\bottomrule
\caption{Підсумкова статистика (інстанс c101)}
\end{longtable}

\begin{longtable}{lcccccc}
\toprule
Алгоритм & min cost & avg cost & max cost & min dist & avg dist & max dist\\ 
\midrule
ACO  & 2423103.44  & 2469665.94  & 2500398.61 & 1901.07  & 1941.11  & 1963.06\\
GA   & 69401.78   & 118311.21  & 161616.55 & 2795.04  & 2920.85 & 3073.03\\
Tabu & 1919.95   & 1946.90 & 1967.86 & 1919.95  & 1946.90 & 1967.86\\
\bottomrule
\caption{Підсумкова статистика (інстанс r101)}
\end{longtable}

\begin{longtable}{lcccccc}
\toprule
Алгоритм & min cost & avg cost & max cost & min dist & avg dist & max dist\\ 
\midrule
ACO  & 27213089.16  & 28124545.31  & 29333216.11 & 1091.34  & 1179.79   & 1345.40\\
GA   & 7935999.73    & 10427469.54  & 13690801.96 & 3066.02   & 3179.34  & 3311.68\\
Tabu & 1879423.91    & 2734560.07  & 4135349.99 & 1500.87  & 1738.27  & 2080.15\\
\bottomrule
\caption{Підсумкова статистика (інстанс c201)}
\end{longtable}

\begin{longtable}{lcccccc}
\toprule
Алгоритм & min cost & avg cost & max cost & min dist & avg dist & max dist\\ 
\midrule
ACO  & 18066407.52  & 19122801.23  & 19818220.15 & 1276.07  & 1402.66    & 1557.06\\
GA   & 904598.66     & 1556908.51   & 1970099.22 & 2724.61   & 2854.76   & 2984.61\\
Tabu & 802890.86     & 1051323.29  & 1294973.43 & 1932.53  & 2229.28   & 2462.77\\
\bottomrule
\caption{Підсумкова статистика (інстанс r201)}
\end{longtable}

\begin{longtable}{lcccccc}
\toprule
Алгоритм & min cost & avg cost & max cost & min dist & avg dist & max dist\\ 
\midrule
ACO  & 1883688.28    & 1931197.18  & 1978589.90 & 2057.32    & 2220.33  & 2599.46\\
GA   & 48706.48     & 60048.44   & 74358.90 & 3356.56   & 3414.44    & 3508.37\\
Tabu & 1949.75      & 2052.95   & 2139.47 & 1949.75   & 2052.95    & 2139.47\\
\bottomrule
\caption{Підсумкова статистика (інстанс rc101)}
\end{longtable}

\begin{longtable}{lcccccc}
\toprule
Алгоритм & min cost & avg cost & max cost & min dist & avg dist & max dist\\ 
\midrule
ACO  & 16937688.93    & 17538615.63   & 18146784.69 & 1124.59    & 1221.99   & 1369.29\\
GA   & 826397.00      & 1140599.06    & 1724271.78 & 3534.04    & 3721.18     & 4147.99\\
Tabu & 411164.46       & 598196.68    & 977916.01 & 2731.54    & 2829.81     & 2966.60\\
\bottomrule
\caption{Підсумкова статистика (інстанс rc201)}
\end{longtable}

\section{Деякі графіки збіжності та маршрутів(інстанс r101)}
\begin{figure}[ht]
\centering
\fbox{\includegraphics[width=.9\linewidth]{D:/University/semestr_6/IntelligentComputing/RR/imgs/aco_r101_convergence.png}}
\caption{Графік збіжності 5 запусків ACO}
\end{figure}
\begin{figure}[ht]
\centering
\fbox{\includegraphics[width=.9\linewidth]{D:/University/semestr_6/IntelligentComputing/RR/imgs/aco_r101_route.png}}
\caption{Графік найкращого маршруту ACO}
\end{figure}

\begin{figure}[ht]
\centering
\fbox{\includegraphics[width=.9\linewidth]{D:/University/semestr_6/IntelligentComputing/RR/imgs/gen_r101_convergence.png}}
\caption{Графік збіжності 5 запусків GA}
\end{figure}
\begin{figure}[ht]
\centering
\fbox{\includegraphics[width=.9\linewidth]{D:/University/semestr_6/IntelligentComputing/RR/imgs/gen_r101_route.png}}
\caption{Графік найкращого маршруту GA}
\end{figure}

\begin{figure}[ht]
\centering
\fbox{\includegraphics[width=.9\linewidth]{D:/University/semestr_6/IntelligentComputing/RR/imgs/tabu_r101_convergence.png}}
\caption{Графік збіжності 5 запусків Tabu}
\end{figure}
\begin{figure}[ht]
\centering
\fbox{\includegraphics[width=.9\linewidth]{D:/University/semestr_6/IntelligentComputing/RR/imgs/tabu_r101_route.png}}
\caption{Графік найкращого маршруту Tabu}
\end{figure}

\clearpage
\newpage

\section{Висновки}
\begin{enumerate}
  \item \textbf{Поставлена мета досягнута.}  
        Реалізовано три евристики — мурашиний Ant Colony Optimization (ACO),
        генетичний GA та траєкторний Tabu Search — з єдиною цільовою
        функцією та модулем логування; створено скрипт, що
        автоматично запускає алгоритми, зберігає історію
        $(\text{iter},t,\text{cost},\text{dist})$ і будує графіки
        збіжності.

  \item \textbf{Фізибельність і якість рішень.}  
        За 60\,с \emph{лише Tabu Search} у деяких інстансів знаходить
        повністю фізибельні тури (\(cost = dist\)); ACO та GA залишають
        штрафи.  
        Tabu Search демонструє найменший середній \emph{cost}
        (≈ 80 \% кращий за GA і ≈ 90 \% — за ACO).  
        ACO дає найкоротшу відстань \(D\), але накопичує великі штрафи;
        GA без локального пошуку суттєво відстає за обома метриками.

  \item \textbf{Напрями покращення.}  
        \begin{itemize}
          \item Додавання локального пошуку (2-opt/relocate) до GA та
                ACO має підвищити якість рішень.
          \item Перехід від звичайної ACS до MACS-VRPTW для мурашиного
                алгоритму очікувано дасть додатковий приріст.
          \item Подальші дослідження доцільно спрямувати на
                автоматичний тюнінг параметрів
                ($\alpha,\beta,\rho$ для ACO; \textit{tenure} — для Tabu)
                та паралельні GPU/мультипроцесорні реалізації.
        \end{itemize}
\end{enumerate}

\section*{Список джерел}
\begin{enumerate}
  \item Dorigo M., Gambardella L.M. Ant Colony System: A Cooperative Learning Approach. \emph{IEEE Trans. Evol. Comput.} 1997.
  \item Solomon M.M. Algorithms for the VRPTW. \emph{Operations Research}. 1987.
  \item Taillard E., Gambardella L., Agazzi G. MACS-VRPTW. 1999.
  \item Laporte G. et al. Handbook of VRP. Springer, 2014.
  \item https://people.idsia.ch/~luca/macs-vrptw/problems/welcome.htm
\end{enumerate}

\end{document}
